\documentclass[11pt]{article}

\usepackage{amsmath,amssymb}
\usepackage{xurl}
\usepackage[hidelinks]{hyperref}
\setlength{\emergencystretch}{2em}

% Shared commands for notes in docs/paper-gaps/.
% Notation follows the LDT paper (references/ldt-paper/).

\newcommand{\Lean}{\textsc{Lean}}
\newcommand{\leanid}[1]{\textsf{#1}}
\newcommand{\ghissue}[1]{\href{https://github.com/LionSR/MIPStarRE/issues/#1}{GitHub issue \##1}}
\newcommand{\ghpr}[1]{\href{https://github.com/LionSR/MIPStarRE/pull/#1}{GitHub PR \##1}}

% --- Paper-compatible notation ---
\newcommand{\F}{\mathbb{F}}
\newcommand{\N}{\mathbb{N}}
\newcommand{\C}{\mathbb{C}}
\newcommand{\tr}{\operatorname{tr}}
\newcommand{\ot}{\otimes}
\newcommand{\E}{\mathop{\bf E\/}}
\newcommand{\bx}{\mathbf{x}}
\newcommand{\by}{\mathbf{y}}
\newcommand{\bu}{\mathbf{u}}
\newcommand{\bv}{\mathbf{v}}

% Approximation relation (paper uses \approx_\eps and \simeq_\zeta)
\newcommand{\approxeps}{\approx_{\varepsilon}}
\newcommand{\simgeq}{\simeq_{\zeta}}

% Polynomial spaces (paper uses \polyfunc, \polysub, \polymeas)
\newcommand{\polyfunc}[3]{\operatorname{Poly}(#1,#2,#3)}
\newcommand{\polysub}[3]{\operatorname{PolySub}(#1,#2,#3)}
\newcommand{\polymeas}[3]{\operatorname{PolyMeas}(#1,#2,#3)}


\title{Review: How the LDT Paper Uses the Self-Improvement SDP}
\author{}
\date{2026-05-05}

\newcommand{\Tr}{\operatorname{Tr}}
\newcommand{\cP}{\mathcal P}

\begin{document}
\maketitle

\section{The SDP in the source}

We examine the only semidefinite program which the low individual degree test
paper uses as an analytic ingredient.  The introduction says that the
self-improvement lemma combines a semidefinite program with the expansion and
Schwartz--Zippel facts used in the classical BFL analysis
\cite[Introduction]{Ji2020LowIndividualDegree}.%
\footnote{The overview sentence is in
\path{references/ldt-paper/introduction.tex}, lines~508--511.  The SDP itself
is in \path{references/ldt-paper/self_improvement.tex}, lines~62--88; the
Slater and complementary-slackness paragraph is in lines~168--190.}
The program is not a separate part of the induction loop.  It is the selection
mechanism inside the self-improvement lemma.

Let $A=\{A^u_a\}$ be the point measurement in the fixed symmetric strategy, and
let
\[
  A_g = \E_{\bu\sim\F_q^m} A^{\bu}_{g(\bu)}
  \qquad (g\in \cP),
  \tag{\texttt{self\_improvement.tex:66--68}}
\]
where $\cP=\polyfunc{m}{q}{d}$ is the finite set of low individual degree
polynomials.  The paper's primal and dual programs are
\begin{align}
  \sup_{\{T_g\}_{g\in\cP}}
    &\quad \sum_{g\in\cP} \Tr(T_g A_g) \\
  \text{subject to}
    &\quad T_g\ge 0 \qquad(g\in\cP), \notag\\
    &\quad \sum_{g\in\cP} T_g\le I, \notag
  \tag{\texttt{eq:primal-objective}}
\end{align}
\begin{align}
  \inf_Z
    &\quad \Tr(Z) \\
  \text{subject to}
    &\quad Z\ge A_g \qquad(g\in\cP). \notag
  \tag{\texttt{eq:dual-objective}}
\end{align}
The primal variable is therefore a submeasurement.  The SDP lemma then asserts
that an optimal pair may be chosen so that
\[
  \sum_{g\in\cP} T_g = I,
  \qquad
  T_g Z = T_g A_g \quad(g\in\cP).
  \tag{\texttt{lem:sdp}, \texttt{eq:slater}}
\]
The equality $\sum_g T_g=I$ is a conclusion about the selected optimum, not the
primal feasibility constraint.

\section{What the selected optimal pair does}

After the SDP lemma, the proof of self-improvement fixes such an optimal pair
$(T,Z)$ and defines the sandwiched family
\[
  H^u_h = A^u_{h(u)} T_h A^u_{h(u)},
  \qquad
  H_h = \E_{\bu} H^{\bu}_h.
  \tag{\texttt{self\_improvement.tex:203--220}}
\]
The downstream proof uses only a small list of mathematical consequences of the
SDP stage.

First, $T$ is used as a submeasurement in positivity and Cauchy--Schwarz
bounds.  The inequalities $T_g\ge0$ and $\sum_g T_g\le I$ imply that each
$H^u$ and the average $H$ are submeasurements.  The stronger saturated equality
$\sum_gT_g=I$ is used in a few scalar normalizations, but those uses occur after
the SDP argument has produced the saturated optimum.

Second, complementary slackness is used to replace $A_g$ by $Z$ under the
left multiplication by $T_g$.  In the paper this is the displayed identity
\[
  T_g \Big(\E_{\bu} A^{\bu}_{g(\bu)}\Big) = T_g Z.
\]
It appears explicitly in the completeness proof, where
\[
  \sum_h \bra{\psi} (T_h A_h)\otimes I\ket{\psi}
  = \sum_h \bra{\psi} (T_h Z)\otimes I\ket{\psi}
  = \bra{\psi} Z\otimes I\ket{\psi},
  \tag{\texttt{self\_improvement.tex:397--400}}
\]
and again in the strong self-consistency proof, where the same swap turns a
sum over $T_h A_h$ into a term involving $Z$.%
\footnote{The second swap is in
\path{references/ldt-paper/self_improvement.tex}, lines~580--584.}
Without this complementary-slackness identity, the proof does not obtain the
paper's comparison between the mass of $H$ and the dual certificate $Z$.

Third, dual feasibility $Z\ge A_g$ supplies the boundedness certificate.  In
the non-projective helper it is part of the output, and after projectivization
the final boundedness statement has the same form: $Z$ is positive semidefinite,
$Z\ge A_h$, and the residual involving $Z\otimes(I-H)$ is small.  The induction
section then carries these $Z$-certificates as hypotheses for later pasting and
restriction steps rather than solving a new SDP at each use.%
\footnote{The boundedness outputs are restated in
\path{references/ldt-paper/inductive_step.tex}, lines~277--284 and 315--322;
the proof produces the per-slice certificates in lines~477--484 and averages
them in lines~542--549.}
Thus the SDP has one analytic role: produce a corrected polynomial
submeasurement together with a dual operator that can witness boundedness.

\section{The submeasurement boundary}

The distinction between the primal feasibility constraint and the saturated
optimal conclusion is mathematically necessary.  If $M=|\cP|$, the strict
Slater witness displayed in the paper is
\[
  T_g = \frac{1}{2M}I,
  \qquad
  Z=2I.
  \tag{\texttt{self\_improvement.tex:171--176}}
\]
The total primal mass is then $\sum_g T_g=\frac12 I$.  It is strictly feasible
for the paper's submeasurement program, but it would be infeasible for an
SDP whose primal variables were required to satisfy $\sum_gT_g=I$ from the
start.  The saturation $\sum_gT_g=I$ is obtained later from complementary
slackness: in the canonical form, the slack block $X_{M+1,M+1}$ vanishes for
the selected optimal solution.

This is the resolved mismatch recorded separately in the SDP weakening note.%
\footnote{See \path{docs/paper-gaps/issue-196-sdp-weakening.tex}.  The original
formalization issue was \ghissue{196}, and the paper-gap documentation issue was
\ghissue{935}.}
The source argument is coherent provided the primal remains a submeasurement
program until the strong-duality and complementary-slackness step has selected a
saturated optimum.

\section{The current formal statement}

The formal development now mirrors this boundary.  The primal objective,
complementary-slackness equation, strict primal witness, and sandwiched-family
construction take the primal variable to be a submeasurement, not a measurement.%
\footnote{The relevant declarations include
\leanid{sdpStrictPrimalSubMeas},
\leanid{sdpPrimalObjective}, \leanid{sdpComplementarySlacknessEquation}, and
\leanid{averagedSandwichedPolynomialSubMeas} in
\path{MIPStarRE/LDT/SelfImprovement/Defs.lean}.}
The formal strict witness has total $(1/2)I$, while a separate completion wrapper
exists only for downstream code paths that need an ordinary measurement.

The current proof interface deliberately separates two layers.  A reduced
statement records the SDP facts already produced by the current wrapper: a
submeasurement witness whose total is $I$, a positive dual operator $Z$ that
dominates the identity, and dual feasibility $Z\ge A_g$.  A stronger statement
adds the complementary-slackness equations required by the paper-facing helper
argument.%
\footnote{These are
\leanid{SdpStatement} and
\leanid{SdpStatementWithSlackness} in
\path{MIPStarRE/LDT/SelfImprovement/Theorems/Statements.lean}.  The helper
wrapper consuming the stronger statement is
\leanid{selfImprovementHelperWithSlackness} in
\path{MIPStarRE/LDT/SelfImprovement/Theorems/Results/SelfImprovementTop.lean}.}
At the finite matrix layer, the analogous strong-duality output is recorded as a
matrix witness with primal identity total, dual feasibility, equality of primal
and dual objective values, and complementary slackness.%
\footnote{See
\leanid{MatrixSdpOptimalWitness} and
\leanid{MatrixSdpStatementWithSlackness} in
\path{MIPStarRE/LDT/SelfImprovement/MatrixRealization.lean}.}

This separation is faithful to the paper, but it is not yet a complete formal
proof of the paper's SDP lemma.  The analytic convex-optimization theorem which
produces the slackness-carrying optimal pair remains a formalization task.  It
is tracked as \ghissue{1230}.  That task should prove the strong-duality
producer; it should not change the primal feasible set from $\sum_gT_g\le I$ to
$\sum_gT_g=I$.

\section{Conclusion}

The LDT paper uses SDP in one place: the self-improvement lemma.  The SDP chooses
a primal polynomial submeasurement $T$ and a dual positive operator $Z$.  The
rest of the proof uses their consequences, namely submeasurement bounds for
sandwiching, dual feasibility $Z\ge A_g$, the saturated optimum
$\sum_gT_g=I$, and complementary slackness $T_gZ=T_gA_g$.

The only mathematical caveat exposed by the review is the already documented
submeasurement boundary: equality of the primal total is a derived property of a
chosen optimum, not a feasibility hypothesis.  The current formalization records
that boundary correctly and has also separated the reduced SDP wrapper from the
stronger slackness-carrying statement.  The remaining work is to formalize the
strong-duality and complementary-slackness producer itself; this is a missing
formal proof of the paper's SDP lemma, not a new discrepancy in the source
argument.

\bibliographystyle{alpha}
\bibliography{references}

\end{document}
